\section{Theta formulas and the basic companions}\label{sec:theta}

\subsection{The simultaneous shift}
Let $m\ge1$, $r=2m+1$, $a=(a_1,\dots,a_m)\in\Q^m$ and $c\in\Q$. Define
\[
B(a,c)=(a_1,-a_1,\dots,a_m,-a_m,c)^\T,\qquad S_m(a,c;q)=F_{2m+1,B(a,c),0}(q).
\]
The following consequence of the rank reduction of Shi and Wang explains a relation visible in the rank-three and rank-five tables. In rank five it specializes to \cite[(5.45)]{SW26} after replacing $q$ by $q^4$.

\begin{lemma}[Simultaneous shift]\label{lem:simshift}
For every integer $h$,
\begin{equation}\label{eq:simshift}
S_m(a+2h\one,c+h;q)=q^{-h\sum_{j=1}^ma_j-mh^2}\,\frac{(-q^{c+h+1/2};q)_\infty}{(-q^{c+1/2};q)_\infty}\,S_m(a,c;q).
\end{equation}
\end{lemma}

\begin{proof}
For $d\in\frac12\Z$, put
\[
\Theta_a(d;q)=\sum_{\ell\in\Z}q^{(\ell-d)^2-a\ell}.
\]
The change of summation variable $\ell\mapsto\ell+h$ gives
\[
\Theta_{a+2h}(d;q)=q^{-ha-h^2-2hd}\,\Theta_a(d;q).
\]
By \cite[Theorem~1.2, (1.11)]{SW26}, setting $n_0=0$ yields
\begin{multline*}
S_m(a,c;q)=\frac{(-q^{c+1/2};q)_\infty}{(q;q)_\infty^m}\sum_{n_1,\dots,n_m\ge0}\frac{q^{\frac14\sum_{j=1}^m(n_j-n_{j-1})^2}\,q^{\sum_{j=1}^{m-1}(a_{j+1}-a_j)n_j+(c-a_m)n_m}}{\prod_{j=1}^m(q;q)_{n_j}}\\
\times\prod_{j=1}^m\Theta_{a_j}\bigl((n_{j-1}-n_j)/2;q\bigr).
\end{multline*}
Under the simultaneous shift, the second exponential gains $q^{-hn_m}$. The product of theta factors gains $q^{-h\sum_ja_j-mh^2+hn_m}$, since $\sum_j(n_{j-1}-n_j)=-n_m$. These factors combine as asserted. All changes of variables are justified by absolute convergence.
\end{proof}

In particular, let $x_j\in\{0,1\}$ and $k=\sum_jx_j$, and define
\[
B_\pm=(\pm1-x_1,\ x_1\mp1,\ \dots,\ \pm1-x_m,\ x_m\mp1,\ \pm\tfrac12)^\T.
\]
Taking $h=-1$ in \cref{lem:simshift} gives the exact identity
\begin{equation}\label{eq:halfshift}
F_{r,B_-,C+k}(q)=2F_{r,B_+,C}(q)\qquad(C\in\Q).
\end{equation}
Thus modularity in the two half-shift families is equivalent. No ordering of the $x_j$ is needed for this identity. As one consequence, the all-rank modularity of the last-coordinate half-shift \cite{CTW26} also gives the modularity, for every odd $r=2m+1$, of
\[
F_{r,(-2,2,\dots,-2,2,-1/2),\,m+\frac{r(4-r)}{24(r+2)}}.
\]

\subsection{The proposed Wronskians}
For the remainder of the paper, let $m\ge1$, $r=2m+1$, $L=2m+3$ and
\[
\nu=\frac{r(r-1)}2=m(2m+1).
\]
For $0\le k\le m$, take $x_1=\cdots=x_{m-k}=0$ and $x_{m-k+1}=\cdots=x_m=1$. Write $B(k,z)$ and $C(k,z)$ for the vector and scalar of \cref{conj:odd}, and put
\[
S^{(z)}_{m,k}(q)=F_{r,B(k,z),0}(q).
\]
It suffices to consider $z=0,\frac12$, since \eqref{eq:halfshift} determines the series with $z=-\frac12$.

Define increasing tuples by deleting the indicated entry:
\begin{align*}
\rho^{(0)}_k&=(1,2,\dots,m+1)\setminus\{m+1-k\}, & \rho^{(0)}_*&=(1,2,\dots,m),\\
\rho^{(1/2)}_k&=(\tfrac12,\tfrac32,\dots,m+\tfrac12)\setminus\{k+\tfrac12\}, & \rho^{(1/2)}_*&=(\tfrac12,\tfrac32,\dots,m-\tfrac12).
\end{align*}
Deletion preserves the order of the remaining entries. For $v=(v_1,\dots,v_m)$, define
\[
\chi(v)=\prod_{i=1}^mv_i\prod_{1\le i<j\le m}(v_i^2-v_j^2),\qquad \|v\|^2=\sum_{i=1}^mv_i^2.
\]

\begin{conjecture}[Companion Wronskian identity]\label{conj:wron}
For $z\in\{0,\frac12\}$ and $0\le k\le m$, writing $\rho=\rho^{(z)}_k$, one has
\begin{equation}\label{eq:wron}
S^{(z)}_{m,k}(q)=\frac{(-q^{1/2+z};q)_\infty^r}{(q;q)_\infty^\nu}\sum_{v\in\rho+L\Z^m}\frac{\chi(v)}{\chi(\rho^{(z)}_*)}\,q^{\frac{\|v\|^2-\|\rho\|^2}{2L}}.
\end{equation}
\end{conjecture}

The cases $(k,z)=(0,0)$ and $(m,\frac12)$ are the principal and last-coordinate half-shift formulas of Chern, Tran and Wakhare \cite[Theorem~1.4]{CTW26}, which are based on \cite{BW15}. \Cref{thm:basic} proves the cases $k\in\{0,1,m\}$, \cref{thm:k2} adds $k=2$, and \cref{thm:k3to6} adds $k=3,4,5,6$.

\begin{proposition}[Modularity of the proposed right-hand sides]\label{prop:rhs-modular}
Let $R^{(z)}_{m,k}$ denote the right-hand side of \eqref{eq:wron}. Then $q^{C(k,z)}R^{(z)}_{m,k}(q)$ is a modular function on a congruence subgroup. Consequently, \cref{conj:wron} implies the sufficiency direction of \cref{conj:odd}.
\end{proposition}

\begin{proof}
For rational $a$, define the unary theta series
\[
\theta_{a,L}(q)=\sum_{v\in a+L\Z}v\,q^{v^2/(2L)},\qquad \mathcal D=q\frac{d}{dq}.
\]
Expanding the determinant and using the Vandermonde identity gives
\[
\det\bigl(\mathcal D^{i-1}\theta_{\rho_j,L}\bigr)_{i,j=1}^m=\frac{(-1)^{\binom m2}}{(2L)^{\binom m2}}\sum_{v\in\rho+L\Z^m}\chi(v)\,q^{\|v\|^2/(2L)}.
\]
The theta series $\theta_{\rho_j,L}$ have weight $3/2$ on a common congruence subgroup. Their Wronskian has weight $3m/2+m(m-1)=\nu/2$: the lower-order derivative terms in its transformation law form a triangular matrix and do not alter the determinant. This is the Wronskian argument used in \cite{MW24,CTW26}.

Put
\[
E_0(\tau)=\frac{\eta(\tau)^2}{\eta(\tau/2)\eta(2\tau)},\qquad E_{1/2}(\tau)=\frac{\eta(2\tau)}{\eta(\tau)}.
\]
The product identities
\[
(-q^{1/2};q)_\infty=q^{1/48}E_0(\tau),\qquad (-q;q)_\infty=q^{-1/24}E_{1/2}(\tau)
\]
and direct evaluation of the sums of squares give
\[
C(k,0)=\frac{\|\rho^{(0)}_k\|^2}{2L}-\frac{r^2}{48},\qquad
C(k,\tfrac12)=\frac{\|\rho^{(1/2)}_k\|^2}{2L}+\frac{r(3-r)}{48}.
\]
It follows that
\[
q^{C(k,z)}R^{(z)}_{m,k}(q)=\frac{E_z(\tau)^r}{\eta(\tau)^\nu\,\chi(\rho^{(z)}_*)}\sum_{v\in\rho^{(z)}_k+L\Z^m}\chi(v)\,q^{\|v\|^2/(2L)}.
\]
The weights cancel. Theta series and eta quotients have finite multiplier systems on congruence subgroups and are meromorphic at the cusps, so restriction to a common congruence subgroup gives the assertion.
\end{proof}

\subsection{Leading coefficients}
The constants in \eqref{eq:wron} are forced by the Nahm sums; they are not fitted separately in each rank.

\begin{proposition}[Leading coefficients]\label{prop:leading}
The series $S^{(0)}_{m,k}$ and $S^{(1/2)}_{m,k}$ have no negative powers of $q$. Their constant terms are, respectively,
\[
d^{(0)}_{m,k}=\binom{2m}{k}-\binom{2m}{k-2},\qquad d^{(1/2)}_{m,k}=\binom{2m+1}{m-k}.
\]
These equal $\chi(\rho^{(z)}_k)/\chi(\rho^{(z)}_*)$, the constant terms of the proposed right-hand sides.
\end{proposition}

\begin{proof}
Put $n_0=0$ and $a_i=n_i-n_{i-1}$, so that $n_j=\sum_{i=1}^ja_i\ge0$. If $E$ denotes the exponent in the unnormalized Nahm summand, then
\[
E=\frac12\sum_{i=1}^ra_i^2+zn_r-\sum_{j=1}^m(2z-x_j)a_{2j}.
\]
For $z=0$ this can be written as
\[
E=\frac{n_r}2+\sum_{i\notin I}\frac{a_i(a_i-1)}2+\sum_{i\in I}\frac{a_i(a_i+1)}2,\qquad I=\{2m-2k+2,\,2m-2k+4,\,\dots,\,2m\}.
\]
For $z=\frac12$ it becomes
\[
E=\sum_{i\in J}\frac{a_i(a_i-1)}2+\sum_{i\notin J}\frac{a_i(a_i+1)}2,\qquad J=\{2,4,\dots,2m-2k\}.
\]
All terms are nonnegative. Since each denominator has constant term one, the constant terms count the paths for which every displayed term vanishes.

We evaluate these path counts. On the basis $e_h$ indexed by integers $h\ge0$, let $Ue_h=e_{h+1}$, let $Ve_0=0$ and $Ve_h=e_{h-1}$ for $h\ge1$. Put $U_+=\Id+U$, $U_-=\Id+V$ and $J_0=U_-U_+=2\Id+U+V$, where $\Id$ is the identity operator. Let $[e_h]f$ denote the coefficient of $e_h$ in $f$, and let $\lambda(f)$ be the sum of the coefficients of $f$. The two zero-energy path counts are
\[
[e_0]\,J_0^kU_+^{2(m-k)}e_0\qquad\text{and}\qquad \lambda\bigl(U_-^{2k}J_0^{m-k}e_0\bigr).
\]
Reflection at height $-1$ gives
\[
[e_a]\,J_0^je_0=\Delta_j(a):=\binom{2j}{j-a}-\binom{2j}{j-a-2}.
\]
Here and below $\binom Nb=0$ unless $0\le b\le N$. By the symmetry of $J_0$ and Vandermonde convolution, the first count is
\[
\sum_{a\ge0}\binom{2(m-k)}a\Delta_k(a)=\binom{2m}k-\binom{2m}{k-2}.
\]
Writing $s=m-k$, the second count is
\[
\sum_{a\ge0}\Delta_s(a)\sum_{d=0}^a\binom{2k}d=\sum_{d\ge0}\binom{2k}d\left(\binom{2s}{s-d}+\binom{2s}{s-d-1}\right)=\binom{2m+1}s.
\]
Finally, direct cancellation in the finite products defining $\chi$ gives the same binomial formulas for the two ratios in the statement. Every coordinate of $\rho^{(z)}_k$ lies strictly between $0$ and $L/2$, so this tuple is the unique point of its lattice coset with minimal norm. Thus its contribution is exactly the constant term of the right-hand side of \eqref{eq:wron}.
\end{proof}

\subsection{The basic companion theorem}

\begin{theorem}[Basic companions]\label{thm:basic}
For every $m\ge1$, identity \eqref{eq:wron} holds for $k\in\{0,1,m\}$ and $z\in\{0,\frac12\}$. Consequently, $F_{2m+1,B(k,z),C(k,z)}$ is a modular function for
\[
k\in\{0,1,m\},\qquad z\in\{-\tfrac12,0,\tfrac12\}.
\]
For $m\ge2$ these are nine distinct pairs.
\end{theorem}

The proof uses two identities of Bartlett and Warnaar \cite[Theorems~3.7 and 5.2]{BW15}; we give the intervening determinant calculation and the specialization factors explicitly.

\subsection{Hall--Littlewood boundary formulas}
Throughout the boundary calculations, $w=q^{1/2}$, $L=2m+3$ and $\nu=m(2m+1)$. From here to the end of \cref{sec:cert}, $x=(x_1,\dots,x_m)$ denotes a vector of Hall--Littlewood variables; the Boolean coordinates of \cref{conj:odd} do not appear in these sections. A partition $\lambda$ is a finite weakly decreasing sequence of positive integers; $|\lambda|$ is its sum, $\ell(\lambda)$ its length, and $m_j(\lambda)$ the multiplicity of $j$. Define
\begin{equation}\label{eq:Hdef}
\HL_m(t;x;q)=\sum_{\lambda:\ \lambda_1\le2}q^{|\lambda|/2}\,t^{m_1(\lambda)}\,P'_\lambda(x_1,x_1^{-1},\dots,x_m,x_m^{-1};q),
\end{equation}
including the empty partition, whose summand is $1$. Here $P'_\lambda(X;q)=P_\lambda[X/(1-q);q]$ is the modified Hall--Littlewood symmetric function: $X/(1-q)$ replaces each letter $X_i$ by $X_i,X_iq,X_iq^2,\dots$. The ordinary Hall--Littlewood functions are determined, for $N\ge\ell(\lambda)$, by
\[
P_\lambda(X_1,\dots,X_N;q)=\frac1{v_\lambda(q)}\sum_{\sigma\in\mathfrak S_N}\sigma\Bigl(X_1^{\lambda_1}\cdots X_N^{\lambda_N}\prod_{i<j}\frac{X_i-qX_j}{X_i-X_j}\Bigr),\qquad v_\lambda(q)=\prod_{j\ge0}\frac{(q;q)_{m_j(\lambda)}}{(1-q)^{m_j(\lambda)}}.
\]
Zero parts are appended to length $N$, so $m_0(\lambda)=N-\ell(\lambda)$, and $\mathfrak S_N$ acts by permuting the variables. Stable specialization defines the symmetric function used in \eqref{eq:Hdef}.

For generic $x=(x_1,\dots,x_m)$ put
\begin{align*}
D_m(x;q)&=(q;q)_\infty^m\prod_{i=1}^m(qx_i^2,qx_i^{-2};q)_\infty\prod_{i<j}(qx_ix_j,qx_i/x_j,qx_j/x_i,q/(x_ix_j);q)_\infty,\\
P_t(x)&=\frac{\prod_i(-wtx_i,-wt/x_i;q)_\infty}{D_m(x;q)},\\
\delta_C(x)&=\det\bigl(x_i^j-x_i^{-j}\bigr)_{i,j=1}^m,\qquad \delta_B(x)=\det\bigl(x_i^{j-1/2}-x_i^{-j+1/2}\bigr)_{i,j=1}^m.
\end{align*}
For the rational functions and half-integral Laurent polynomials below, at values without poles on $xq^\Z$, set
\[
\mu_x[f]=\sum_{\ell\in\Z}q^{L\ell^2/2}x^{L\ell}f(xq^\ell),\qquad \Phi_a(x)=\mu_x[y^a-y^{-a}],
\]
and
\[
\Psi_j(x)=\mu_x\left[\frac{y^j-y^{-j}}{1+w^2+w(y+y^{-1})}\right].
\]
The letter $y$ denotes the argument of $f$ inside $\mu_x$. Powers are taken with a fixed branch of the logarithm on a neighbourhood of the relevant points: the evaluations below are first made for $0<q<1$ and positive variables, and then continued analytically. Gaussian decay gives normal convergence away from the poles.

\begin{lemma}[Boundary specialization]\label{lem:boundary}
For rational $v,c$, at points where the expressions are defined,
\begin{equation}\label{eq:boundary}
S_m(-2v\one,c;q)=(-q^{c+1/2};q)_\infty\,\HL_m(q^{c+v};q^v\one;q).
\end{equation}
For generic $x$ one has
\begin{align}
\HL_m(1;x;q)&=P_1(x)\frac{\det(\Phi_j(x_i))_{i,j=1}^m}{\delta_C(x)},\label{eq:H1}\\
\HL_m(q;x;q)&=P_1(x)\frac{\det(\Psi_j(x_i))_{i,j=1}^m}{\delta_C(x)},\label{eq:Hq}\\
\HL_m(w;x;q)&=P_w(x)\frac{\det(\Phi_{j-1/2}(x_i))_{i,j=1}^m}{\delta_B(x)}.\label{eq:Hw}
\end{align}
Specializations at coincident variables mean the removable limits.
\end{lemma}

\begin{proof}
Sum over the last coordinate of the Nahm sum using Euler's identity
\[
\sum_{n\ge0}\frac{q^{n(n-1)/2}z^n}{(q;q)_n}=(-z;q)_\infty.
\]
The resulting factor is $(-q^{c+1/2-h};q)_\infty$, where $h$ is the remaining last coordinate. Extract $(-q^{c+1/2};q)_\infty$ and reverse the order of the remaining $2m$ coordinates. In the notation of \cite[(3.21)]{BW15}, the series becomes $F_{1,2m}(q^v,0,q^{c+v};q)$. Their Theorem~3.7 holds for these parameters and identifies this expression with the series on the right of \eqref{eq:boundary}.

For the determinant formulas, take the level parameter equal to $1$ in \cite[(5.6a)]{BW15}, let their parameter $b$ tend to infinity, and set their parameter $c$ equal to $-w/t$. Their equation (5.7) then gives \eqref{eq:Hdef}. After rewriting the type-$C$ denominator as $\delta_C$, the summand has one-variable multiplier
\[
q^{L\ell^2/2}x^{L\ell}t^\ell\,\frac{(-wx/t;q)_\ell}{(-wtx;q)_\ell},
\]
where $(a;q)_\ell=(a;q)_\infty/(aq^\ell;q)_\infty$ for $\ell\in\Z$. The quotient is $1$ at $t=1$, which proves \eqref{eq:H1}. At $t=q$, combine its two telescoping factors with $P_q/P_1$; the resulting multiplier is
\[
\frac{y/w}{(1+y/w)(1+wy)}=\frac1{1+w^2+w(y+y^{-1})},\qquad y=xq^\ell.
\]
This proves \eqref{eq:Hq}. At $t=w$ the quotient is $(1+x)/(1+xq^\ell)$. Using
\[
\delta_C(y)=\prod_i\bigl(y_i^{1/2}+y_i^{-1/2}\bigr)\,\delta_B(y)
\]
to cancel this quotient and $w^\ell$ row by row gives \eqref{eq:Hw}. These calculations hold initially where the cited bilateral series converge without poles; analytic continuation gives the stated specializations. The limits needed below are computed explicitly.
\end{proof}

\subsection{A finite tridiagonal calculation}
For an increasing tuple $a=(a_1,\dots,a_m)$ define
\[
J_a(x)=P_1(x)\frac{\det(\Phi_{a_j}(x_i))_{i,j=1}^m}{\delta_C(x)},\qquad J_k(x)=J_{\rho^{(0)}_k}(x).
\]

\begin{lemma}\label{lem:tridiag}
For generic $x$,
\begin{equation}\label{eq:tridiag}
\HL_m(1;x;q)-\HL_m(q;x;q)=wJ_1(x).
\end{equation}
\end{lemma}

\begin{proof}
Put $d=m+1$, so that $L=2d+1$. Reindexing the bilateral sum gives
\[
\mu_x[f(y)]=q^{L/2}\mu_x[y^Lf(qy)].
\]
Together with $q=w^2$ this implies
\[
\Psi_{d+1}+w\Psi_d=\mu_x\left[\frac{y^{d+1}}{1+wy}-\frac{y^{-d-1}}{1+w/y}\right]=0,
\]
since the displayed reindexing transforms the second rational term into the first. Multiplying by the denominator in the definition of $\Psi_j$ also gives
\[
\Phi_j=(1+w^2)\Psi_j+w(\Psi_{j-1}+\Psi_{j+1}),\qquad \Psi_0=0.
\]
Thus $(\Phi_1,\dots,\Phi_d)^\T=U_d(\Psi_1,\dots,\Psi_d)^\T$, where $U_d$ is tridiagonal with entries $w$ above and below the diagonal and diagonal $1+w^2,\dots,1+w^2,1$. Direct multiplication gives
\[
\det U_d=1,\qquad (U_d^{-1})_{ij}=(-w)^{|i-j|}\sum_{s=0}^{\min(i,j)-1}w^{2s}.
\]
By the complementary-minor formula, the minor of $U_d^{-1}$ formed by its first $m$ rows with column $t$ deleted is $(-1)^{d+t}(U_d)_{t,d}$. It equals $1$ for $t=d$, $-w$ for $t=d-1$, and $0$ otherwise. The Cauchy--Binet formula therefore yields
\[
\det(\Psi_j(x_i))_{i,j=1}^m=\det(\Phi_1,\dots,\Phi_m)-w\det(\Phi_1,\dots,\Phi_{m-1},\Phi_{m+1}),
\]
where each determinant evaluates its columns at $x_1,\dots,x_m$. Equations \eqref{eq:H1} and \eqref{eq:Hq} now prove the lemma.
\end{proof}

\subsection{Confluent limits and the affine translation}
For an increasing tuple $a$ put
\[
N_a(q)=\sum_{v\in a+L\Z^m}\chi(v)\,q^{(\|v\|^2-\|a\|^2)/(2L)},\qquad A_0=\frac{(-w;q)_\infty^{2m}}{(q;q)_\infty^\nu},\qquad A_{1/2}=\frac{(-q;q)_\infty^{2m}}{(q;q)_\infty^\nu}.
\]

\begin{lemma}\label{lem:confluent}
For every $0\le k\le m$,
\begin{align}
J_k(\one)&=A_0\,\frac{N_{\rho^{(0)}_k}}{\chi(\rho^{(0)}_*)},\label{eq:conf0}\\
J_k(w\one)&=q^{-k/2}A_{1/2}\,\frac{N_{\rho^{(1/2)}_k}}{\chi(\rho^{(1/2)}_*)}.\label{eq:confhalf}
\end{align}
Moreover,
\begin{align}
\HL_m(w;\one;q)&=A_{1/2}\,\frac{N_{\rho^{(1/2)}_m}}{\chi(\rho^{(1/2)}_*)},\label{eq:Hw1}\\
\HL_m(w;w\one;q)&=A_0\,\frac{N_{\rho^{(0)}_m}}{\chi(\rho^{(0)}_*)}.\label{eq:Hww}
\end{align}
\end{lemma}

\begin{proof}
Set $x_i=e^{\varepsilon t_i}$ with distinct generic $t_i$, and expand each column in odd powers of $\varepsilon t_i$. Cancellation of the leading Vandermonde determinants gives
\[
\lim_{x\to\one}\frac{\det(\Phi_{a_j}(x_i))}{\delta_C(x)}=\frac{N_a}{\chi(\rho^{(0)}_*)}.
\]
The same calculation with $\delta_B$ gives the denominator $\chi(\rho^{(1/2)}_*)$. Normal convergence justifies the Taylor expansions. Evaluating the nonsingular infinite products proves \eqref{eq:conf0} and \eqref{eq:Hw1}.

Reindexing the theta series gives the exact affine translation
\begin{equation}\label{eq:affine}
\Phi_a(wu)=-q^{-a/2}u^{-L/2}\,\Phi_{L/2-a}(u).
\end{equation}
Write $h=\binom m2$ and $s=m+h$. Product formulas for the two finite determinants give, with $Z(u)=\prod_i(1-u_i^{-2})\prod_{i<j}(1-(u_iu_j)^{-1})$,
\begin{align*}
\lim_{u\to\one}\frac{D_m(wu;q)}{Z(u)}&=\frac{(q;q)_\infty^\nu}{(1-q)^s},\\
\lim_{u\to\one}\frac{\delta_C(wu)Z(u)}{\delta_C(u)}&=(w-w^{-1})^s,\\
\lim_{u\to\one}\frac{\delta_B(wu)Z(u)}{\delta_C(u)}&=(w^{1/2}-w^{-1/2})^m(w-w^{-1})^h.
\end{align*}
For example, $\delta_C(u)=\prod_i(u_i-u_i^{-1})\prod_{i<j}(u_j-u_i)(1-(u_iu_j)^{-1})$, which makes the cancellations explicit; the analogous formula for $\delta_B$ replaces each first factor by $u_i^{1/2}-u_i^{-1/2}$. Consequently,
\begin{align*}
\lim_{u\to\one}\frac{P_1(wu)\,\delta_C(u)}{\delta_C(wu)}&=(-1)^s2^mq^{s/2}A_{1/2},\\
\lim_{u\to\one}\frac{P_w(wu)\,\delta_C(u)}{\delta_B(wu)}&=(-1)^sq^{m/4+h/2}A_0.
\end{align*}
In \eqref{eq:affine}, reversing the order of the transformed columns contributes $(-1)^h$, and the $m$ minus signs contribute $(-1)^m$; together these cancel $(-1)^s$. For $a=\rho^{(0)}_k$ the transformed increasing tuple is $\rho^{(1/2)}_k$, and $\sum_ja_j=s+k$. Since $\chi(\rho^{(1/2)}_*)=2^{-m}\chi(\rho^{(0)}_*)$, the remaining scalar is exactly $q^{-k/2}$, as in \eqref{eq:confhalf}. For $a=\rho^{(1/2)}_*$ the transformed tuple is $(2,3,\dots,m+1)=\rho^{(0)}_m$, and $\sum_ja_j=m^2/2$; the scalar cancels because $m/4+h/2=m^2/4$. Applying \eqref{eq:Hw} proves \eqref{eq:Hww}. These product cancellations also show that all the limits used exist.
\end{proof}

\subsection{Contiguity and completion of the proof}
For brevity write $F_B=F_{2m+1,B,0}$. Subtracting adjacent summands in the defining series gives
\begin{equation}\label{eq:contig}
F_B-F_{B+e_i}=q^{(T_{2m+1})_{ii}/2+b_i}\,F_{B+T_{2m+1}e_i}.
\end{equation}
Take $B=B(a\one,c)$ and $E=e_{2m}-e_{2m-1}$. Apply \eqref{eq:contig} first at the last coordinate and then at coordinate $2m$ of the shifted vector. Since $T_{2m+1}(e_{2m}+e_{2m+1})=E$, the result is
\[
F_{B+E}=q^a\Bigl(q^{-c-1/2}F_B-(1+q^{-c-1/2})F_{B+e_{2m+1}}\Bigr).
\]
With $a=-2v$, $c=-v$ and \cref{lem:boundary}, this becomes
\begin{equation}\label{eq:contig2}
F_{B+E}=q^{-v-1/2}(-q^{1/2-v};q)_\infty\bigl(\HL_m(1;q^v\one;q)-\HL_m(q;q^v\one;q)\bigr).
\end{equation}

\begin{proof}[Proof of \cref{thm:basic}]
For $z=0$, \cref{lem:boundary} and \eqref{eq:contig2} at $v=0$ give
\[
S^{(0)}_{m,0}=(-w;q)_\infty J_0(\one),\qquad S^{(0)}_{m,1}=(-w;q)_\infty J_1(\one),
\]
the second via \cref{lem:tridiag}. For $k=m$, use \eqref{eq:boundary} with $v=\frac12$, $c=0$:
\[
S^{(0)}_{m,m}=(-w;q)_\infty\HL_m(w;w\one;q).
\]
Equations \eqref{eq:conf0} and \eqref{eq:Hww} prove the three identities with $z=0$.

At $v=\frac12$, $c=-\frac12$, identity \eqref{eq:boundary} concerns the negative half-shift with $k=0$. Using \eqref{eq:halfshift} and $(-1;q)_\infty=2(-q;q)_\infty$ we obtain
\[
S^{(1/2)}_{m,0}=(-q;q)_\infty J_0(w\one).
\]
Similarly, \eqref{eq:contig2} at $v=\frac12$ concerns the negative half-shift with $k=1$. Since $S^{(-1/2)}_{m,1}=2q^{-1}S^{(1/2)}_{m,1}$, it gives
\[
S^{(1/2)}_{m,1}=w(-q;q)_\infty J_1(w\one).
\]
Finally, \eqref{eq:boundary} at $v=0$, $c=\frac12$ yields
\[
S^{(1/2)}_{m,m}=(-q;q)_\infty\HL_m(w;\one;q).
\]
Equations \eqref{eq:confhalf} and \eqref{eq:Hw1} give the half-shift identities. \Cref{prop:rhs-modular} supplies modularity with exactly the scalar $C(k,z)$, and \eqref{eq:halfshift} supplies the sector $z=-\frac12$. When $m=1$ some of the indices coincide, consistently with the six rank-three pairs.
\end{proof}

For later use, set
\[
B_v=B(-2v\one,-v)\qquad(v\in\Q).
\]
By \cref{lem:boundary}, every last-coordinate shift of $B_v$ is a Hall--Littlewood boundary specialization.

\subsection{A worked example}\label{subsec:example}
We follow one pair through the paper: rank five ($m=2$, $L=7$, $\nu=10$), index $k=1$ and sector $z=\frac12$. In \cref{conj:odd} this is $x=(0,1)$, so
\[
B(1,\tfrac12)=(1,-1,0,0,\tfrac12)^\T,\qquad C(1,\tfrac12)=-\frac{55}{336}+\frac{5}{14}-\frac12+\frac9{16}=\frac{43}{168}.
\]
In the coordinates of \cref{thm:r5} this point is $u=v=0$, $(x,y)=(0,1)$, $z=\frac12$, and \eqref{eq:r5C} gives the same scalar $(120-168+189-55)/336=43/168$; it is the sixth row of the rank-five table, whose modularity is \cite[(5.34)]{SW26}.

On the companion side, $\rho^{(1/2)}_1=(\frac12,\frac52)$ and $\rho^{(1/2)}_*=(\frac12,\frac32)$, so $\chi(\rho^{(1/2)}_1)=-\frac{15}2$ and $\chi(\rho^{(1/2)}_*)=-\frac32$. Identity \eqref{eq:wron} reads
\begin{align*}
F_{5,(1,-1,0,0,1/2),0}(q)&=\frac{(-q;q)_\infty^5}{(q;q)_\infty^{10}}\sum_{v\in(\frac12,\frac52)+7\Z^2}\frac{v_1v_2(v_1^2-v_2^2)}{-3/2}\,q^{(v_1^2+v_2^2-13/2)/14}\\
&=5+45q+200q^2+715q^3+\cdots,
\end{align*}
and the constant term $5=\binom{5}{1}$ agrees with \cref{prop:leading}. This case is covered by \cref{thm:basic}: \eqref{eq:contig2} at $v=\frac12$ expresses the negative half-shift as a difference of two Hall--Littlewood boundary specializations, \cref{lem:tridiag} turns that difference into the single theta minor $wJ_1$, and \eqref{eq:confhalf} evaluates its confluent limit at $x=w\one$. By \cref{prop:rhs-modular}, $q^{43/168}$ times the right-hand side is a modular function, in agreement with the necessity computation.
